
\section{Element}

To accommodate parameters that may affect the element length, reference geometry, or interpolation data, it is convenient to express the element resisting force on a fixed parent domain \(\hat e\):
\begin{equation}
\bm{p}_a(\ModelState,\HostParameter)
=
\int_{\hat e}
J(\hat{\reflenx},\HostParameter)\,\bm{B}_a(\ModelState,\HostParameter)\,\bm{\sigma}(\ModelState,\HostParameter)\,\dd\hat{\reflenx}
=
\int_{\hat e}
J\,\bm{B}_a\,\bm{A}_*\,\bm{s}\,\dd\hat{\reflenx} .
\label{eq:pa-parent}
\end{equation}
Here \(\ModelState\) denotes the vector of discrete configuration variables in the chosen chart, \(J\) is the parent-to-element Jacobian, and \(\bm{\sigma}=\bm{A}_*\bm{s}\).

Define
\begin{equation}
(\,\cdot\,)_{,\HostParameter} \coloneq \frac{\dd}{\dd \HostParameter}(\,\cdot\,),
\qquad
(\,\cdot\,)_{;\HostParameter}
\coloneq
\left.\frac{\partial}{\partial \HostParameter}(\,\cdot\,)\right|_{\ModelState},
\label{eq:def-direct-h}
\end{equation}
so that \((\,\cdot\,)_{;\HostParameter}\) denotes the explicit parameter derivative at fixed discrete state. Differentiating \eqref{eq:pa-parent} along the equilibrium path gives
\begin{equation}
\bm{p}_{a,\HostParameter}
=
\int_{\hat e}
\left[
D_\chi[\bm{B}_a\bm{A}_*\mid\bm{s}] \cdot \bm{u}_h
+
J\,\bm{B}_a\bm{A}_*\,\bm{s}_{,\HostParameter}
+
\bigl(J\,\bm{B}_a\bm{A}_*\mid\bm{s}\bigr)_{;\HostParameter}
\right]\,\dd\hat{\reflenx} ,
\label{eq:pa-h-raw}
\end{equation}
where \(\bm{u}_h\) is the configuration sensitivity field induced by the nodal sensitivities,
\begin{equation}
\bm{u}_h^h
=
\sum_b \bm{\varphi}_{u b}\,\ModelState_{b,h}.
\label{eq:uh-field}
\end{equation}

The first term in \eqref{eq:pa-h-raw} is the state-induced geometric contribution. 
By definition of the discrete geometric operator,
\begin{equation}
D_\chi[\bm{B}_a\bm{A}_*\mid\bm{s}] \cdot \bm{u}_h
=
\sum_b \bm{G}_{ab}[\bm{s}]\,\ModelState_{b,h}.
\label{eq:geom-state-sens}
\end{equation}
The constitutive term is treated by differentiating the section response. The strain sensitivity admits the decomposition
\begin{equation}
\bm{e}_{,\HostParameter}
=
D\bm{e}\cdot \bm{u}_h + \bm{e}_{;\HostParameter}
=
\sum_b \bm{B}_b^\star \ModelState_{b,h} + \bm{e}_{;\HostParameter},
\label{eq:eh-split}
\end{equation}
and therefore
\begin{equation}
\bm{s}_{,\HostParameter}
=
\bm{C}_s\,\bm{e}_{,\HostParameter} + \bm{s}_{;\HostParameter}
=
\bm{C}_s\sum_b \bm{B}_b^\star \ModelState_{b,h}
+
\bm{C}_s\,\bm{e}_{;\HostParameter}
+
\bm{s}_{;\HostParameter}.
\label{eq:sh-split}
\end{equation}
Substituting \eqref{eq:geom-state-sens} and \eqref{eq:sh-split} into \eqref{eq:pa-h-raw} yields
\begin{equation}
\bm{p}_{a,\HostParameter}
=
\sum_b
\left[
\int_{\hat e}
J\Bigl(
\bm{B}_a\,\bm{C}_\sigma\,\bm{B}_b^\star
+
\bm{G}_{ab}[\bm{s}]
\Bigr)\,\dd\hat{\reflenx}
\right]\ModelState_{b,h}
+
\bm{p}_{a;\HostParameter},
\label{eq:pa-h-final}
\end{equation}
with
\begin{equation}
\bm{p}_{a;\HostParameter}
=
\int_{\hat e}
\left[
\bigl(J\,\bm{B}_a\bm{A}_*\mid\bm{s}\bigr)_{;\HostParameter}
+
J\,\bm{B}_a\bm{A}_*
\Bigl(
\bm{C}_s\,\bm{e}_{;\HostParameter}
+
\bm{s}_{;\HostParameter}
\Bigr)
\right]\,\dd\hat{\reflenx} .
\label{eq:pa-h-direct}
\end{equation}
The first term in \eqref{eq:pa-h-direct} expands to
\begin{equation}
\bigl(J\,\bm{B}_a\bm{A}_*\mid\bm{s}\bigr)_{;\HostParameter}
=
J_{,\HostParameter}\,\bm{B}_a\bm{\sigma}
+
J\,(\bm{B}_a\bm{A}_*)_{;\HostParameter}\,\bm{s}.
\label{eq:pa-h-direct-expand}
\end{equation}

Equation \eqref{eq:pa-h-final} gives the component sensitivity of the resisting force vector. In matrix form,
\begin{equation}
\bm{p}_{,\HostParameter}
=
\bm{k}\,\ModelState_{,\HostParameter}
+
\bm{p}_{;\HostParameter},
\label{eq:p-h-matrix}
\end{equation}
with element tangent
\begin{equation}
\bm{k}_{ab}
=
\int_{\hat e}
J\Bigl(
\bm{B}_a\,\bm{C}_\sigma\,\bm{B}_b^\star
+
\bm{G}_{ab}[\bm{s}]
\Bigr)\,\dd\hat{\reflenx} .
\label{eq:k-ab-sens}
\end{equation}
Thus the force sensitivity splits into a state-induced term, obtained by applying the consistent tangent to the nodal sensitivity vector, and a direct term that collects all explicit parameter dependence of the element geometry, interpolation, and constitutive response.

% For a pure material parameter, such as \(E\) or \(G\), the geometry and interpolation are fixed, so that
% \[
% J_{,\HostParameter}=0,
% \qquad
% (\bm{B}_a\bm{A}_*)_{;\HostParameter}=0,
% \qquad
% \bm{e}_{;\HostParameter}=0,
% \]
% and \eqref{eq:pa-h-direct} reduces to
% \begin{equation}
% \bm{p}_{a;\HostParameter}
% =
% \int_{\hat e}
% J\,\bm{B}_a\bm{A}_*\,\bm{s}_{;\HostParameter}\,\dd\hat{\reflenx}.
% \label{eq:pa-h-material}
% \end{equation}
% If the constitutive law is linear in the form \(\bm{s}=\bm{C}_s(h)\bm{e}\), then
% \begin{equation}
% \bm{s}_{;\HostParameter}=\bm{C}_{s,h}\bm{e},
% \qquad
% \bm{p}_{a;\HostParameter}
% =
% \int_{\hat e}
% J\,\bm{B}_a\bm{A}_*\,\bm{C}_{s,h}\bm{e}\,\dd\hat{\reflenx}.
% \label{eq:pa-h-linear-material}
% \end{equation}

At the assembled level, if the equilibrium residual is written as
\begin{equation}
\bm{g}(\ModelState,\HostParameter)=\bm{p}(\ModelState,\HostParameter)-\bm{f}(\ModelState,\HostParameter),
\label{eq:residual-sens}
\end{equation}
then differentiation gives
\begin{equation}
\bm{g}_{,\HostParameter}
=
\bm{K}\,\ModelState_{,\HostParameter}
+
\bm{p}_{;\HostParameter}
-
\bm{f}_{,\HostParameter},
\label{eq:residual-sens-final}
\end{equation}
where \(\bm{K}\) is the assembled consistent tangent formed from \eqref{eq:k-ab-sens}. 
This is the form required in direct differentiation procedures for nonlinear static and dynamic sensitivity analysis.
